\lyseis{Β2.3}{Μεσοκάθετος ευθυγράμμου τμήματος}
\sol{1} Ελεύθερη κατασκευή. $\mathrm{ΑΜ}=3$ cm και η γωνία $\widehat{\mathrm{ΑΜΚ}}=90^\circ$.
\sol{2} Για κάθε σημείο της μεσοκαθέτου οι δύο αποστάσεις είναι ίσες.
\sol{3} $\mathrm{ΜΒ}=5$.
\sol{4} Πάνω στη μεσοκάθετο του $\mathrm{ΑΒ}$.
\sol{5} Το σημείο τομής των δύο μεσοκαθέτων ισαπέχει από τις τρεις κορυφές.
\sol{6} α) Σωστό. β) Λάθος: πρέπει να είναι και κάθετη. γ) Σωστό. δ) Λάθος.
\sol{7} Τα δύο σημεία της μεσοκαθέτου με $\mathrm{ΜΑ}=\mathrm{ΜΒ}=5$: από το μέσο στα $3$ μ. προς κάθε μεριά (από το τρίγωνο με πλευρές $3$, $4$, $5$: το μισό του $\mathrm{ΑΒ}$ είναι $4$).
\sol{8} Πάνω στη μεσοκάθετο του $\mathrm{ΑΒ}$, δηλαδή στο μέσο ή σε κάποιο άλλο σημείο της.
\sol{9} Φέρνουν τις μεσοκαθέτους δύο ζευγών σπιτιών: το σημείο τομής ισαπέχει από τα τρία.
\sol{10} Το κέντρο είναι το σημείο τομής των μεσοκαθέτων των δύο χορδών.
\sol{11} Το σημείο τομής της $\varepsilon$ με τη μεσοκάθετο του $\mathrm{ΑΒ}$.
\sol{12} Αν η ακτίνα ήταν μικρότερη ή ίση από το μισό του $\mathrm{ΑΒ}$, οι κύκλοι δεν θα τέμνονταν σε δύο σημεία.
\sol{13} Επειδή το $\mathrm{ΑΒ}$ είναι κάθετο στη μεσοκάθετο και η μεσοκάθετος περνά από το μέσο του $\mathrm{ΑΒ}$. Το συμμετρικό του $\mathrm{Α}$ είναι το $\mathrm{Β}$.
\sol{14} Είναι η μεσοκάθετος του $\mathrm{ΑΒ}$, δηλαδή κάθετη στο $\mathrm{ΑΒ}$.
\sol{15} Δύο, πάνω στη μεσοκάθετο, σε απόσταση $4$ cm από το μέσο.
\sol{16} Ελεύθερη κατασκευή.
\sol{17} Ελεύθερη κατασκευή.
